\documentclass[11pt]{article}

\usepackage[margin=1in]{geometry}
\usepackage{amsmath, amssymb, amsthm}
\usepackage{natbib}
\usepackage{hyperref}
\usepackage{booktabs}
\usepackage{enumitem}
\usepackage{algorithm}
\usepackage{algpseudocode}

\setlength{\parskip}{4pt}
\setlength{\parindent}{0pt}

\title{Posterior Summaries for Data Fusion}
\author{Working notes}
\date{\today}

\newcommand{\E}{\mathbb{E}}
\newcommand{\X}{X}
\newcommand{\Y}{Y}
\newcommand{\A}{A}
\newcommand{\Sx}{S}
\newcommand{\U}{U}
\newcommand{\tauf}{\tau}
\newcommand{\cf}{c}
\newcommand{\gf}{g}
\newcommand{\muf}{\mu}
\newcommand{\mzero}{m_{0}}
\newcommand{\Spot}{S}
\newcommand{\RCT}{\mathrm{RCT}}
\newcommand{\OS}{\mathrm{OS}}
\newcommand{\RMST}{\mathrm{RMST}}
\newcommand{\RD}{\mathrm{RD}}

\begin{document}
\maketitle

\section*{Scope}



\section{Posterior projection summaries for interpretable inference}

A practical difficulty of Bayesian-nonparametric HTE models is that the
posterior of $\tauf(\X)$ lives in an effectively infinite-dimensional space:
each MCMC draw is a function, not a vector of coefficients. To produce
inferentially valid summaries on an interpretable scale --- coefficients on
a small set of covariates, subgroup contrasts, tree-style rules --- we adopt
the posterior projection approach of \citet{Woody2020}. The idea is to map
each posterior draw into a simpler function class via a deterministic
projection and to treat the collection of projected summaries as the
posterior over the summary. No model is refit; uncertainty is inherited
exactly.

\subsection{The projection}

Let $\mathcal{G} = \{g_{\gamma}(x) = \phi(x)^{\top}\gamma : \gamma \in
\mathbb{R}^{k}\}$ be a target function class indexed by a coefficient vector
$\gamma$, with basis $\phi : \mathcal{X} \to \mathbb{R}^{k}$. Typical
choices for $\phi$ include main effects of a chosen subset of covariates,
low-order interactions, B-spline bases for continuous covariates, or
indicator dummies for a small set of subgroups.

Given a posterior draw of the target function $\theta^{(r)}(\cdot)$, where
$\theta$ stands for any of $\tauf$, $\cf$, $\Delta_{\RMST}(t^{*}, \cdot)$, or
$\Delta_{\RD}(t, \cdot)$, the projection at iteration $r$ is the solution to
the weighted least squares problem
\begin{equation}
  \gamma^{(r)} = \arg\min_{\gamma \in \mathbb{R}^{k}}
  \sum_{i=1}^{n_{\mathrm{ev}}}
  w_{i}^{(r)}\bigl(\theta^{(r)}(x_{i}) - \phi(x_{i})^{\top}\gamma\bigr)^{2}
  = (\Phi^{\top} W^{(r)} \Phi)^{-1}\Phi^{\top} W^{(r)} \theta^{(r)},
  \label{eq:projection}
\end{equation}
where $\Phi$ is the $n_{\mathrm{ev}} \times k$ design matrix with rows
$\phi(x_{i})^{\top}$, $\theta^{(r)} = (\theta^{(r)}(x_{1}), \ldots,
\theta^{(r)}(x_{n_{\mathrm{ev}}}))^{\top}$, and $W^{(r)} =
\mathrm{diag}(w_{1}^{(r)}, \ldots, w_{n_{\mathrm{ev}}}^{(r)})$ is an
optional weight matrix used either to focus on a target population or to
inject covariate-distribution uncertainty via a Bayesian bootstrap.

The collection $\{\gamma^{(r)}\}_{r=1}^{R}$ is the posterior over the
projection. Posterior means, marginal credible intervals, and probabilities
of the form $\Pr(\gamma_{k} > 0 \mid \text{data})$ are computed as empirical
operations on this collection.

\subsection{Validity}

Because $\gamma^{(r)}$ is a deterministic function of the posterior draw
$\theta^{(r)}$ (and the bootstrap weights $w^{(r)}$, when used), the
collection $\{\gamma^{(r)}\}$ is a sample from the posterior distribution of
the pushforward of $\theta$ under the projection map. No additional
modelling assumption is introduced. The credible intervals derived from
$\{\gamma^{(r)}\}$ inherit the full uncertainty of the original posterior
of $\theta$, in contrast to the common workflow of fitting an interpretable
model to a point summary (e.g.\ the posterior mean) and reporting that
model's standard errors, which double-uses the data and ignores the
original posterior variance.

\subsection{What to project, and on which scale}

In our setting three projections are natural:

\paragraph{(P1) The CATE on the AFT scale: $\tauf(\X)$.}
The cleanest interpretation: $\gamma_{k}$ is the log-time shift per unit of
covariate $k$. Useful for statistical reporting and for diagnostics of the
structure of heterogeneity.

\paragraph{(P2) The survival-scale contrast: $\Delta_{\RMST}(t^{*}, \X)$.}
For clinically meaningful $t^{*}$ (e.g.\ in-trial and well past trial
follow-up), the projection coefficients have the interpretation of
``RMST gain per unit of $x$ at horizon $t^{*}$''. This is the headline
summary for the audience interested in the long-term goal (G2). Reporting
projections at multiple $t^{*}$ traces how the heterogeneity summary moves
into the extrapolation regime.

\paragraph{(P3) The confounding function: $\cf(\X)$.}
A diagnostic projection rather than a clinical summary. The projection of
$\cf(\X)$ describes which covariates carry the residual bias signal in the
OS. If $\cf$ admits a clean low-dimensional projection, the bias structure
is interpretable; if it does not, the OS contribution is intrinsically
complex and the reader should be warned.

\subsection{The headline use: RCT-only vs.\ fusion contrast}

The single most informative deployment of the projection in this paper is
to fit the model under two data configurations --- RCT-only and full
RCT$+$OS fusion --- and to project both posteriors onto the \emph{same}
basis $\phi$ at the \emph{same} evaluation points. The marginal posteriors
of each $\gamma_{k}$ can then be plotted side by side, giving a direct,
scale-matched answer to the question ``how did borrowing change the
heterogeneity story?''. The same construction with the model fit once and
the evaluation points changed (RCT-population vs.\ OS-population
empirical distribution) decomposes ``model effect'' from
``population effect'' in the apparent change of heterogeneity.

\subsection{Bayesian-bootstrap variant}

When the summary target is a functional of the covariate distribution
(e.g.\ a marginal effect, or any quantity defined by averaging over the
empirical distribution), the posterior of the model parameters alone does
not capture sampling uncertainty in that distribution. Following
\citet{Rubin1981} and \citet{Antonelli2019}, we inject covariate-distribution
uncertainty by drawing Dirichlet weights at each MCMC iteration:
\begin{equation}
  w^{(r)} = (w_{1}^{(r)}, \ldots, w_{n_{\mathrm{ev}}}^{(r)}) \sim
  \mathrm{Dirichlet}(1, 1, \ldots, 1),
  \label{eq:bb}
\end{equation}
and use these as the weights in \eqref{eq:projection}. The same Dirichlet
draw is used everywhere it appears in the iteration (e.g.\ when also
computing a marginal RMST as a weighted average over evaluation points),
so that the dependence is preserved.

\subsection{Algorithm}

\begin{algorithm}[h]
\caption{Posterior projection summary}
\label{alg:projection}
\begin{algorithmic}[1]
\Require posterior draws of model parameters $\{\theta_{\text{model}}^{(r)}\}_{r=1}^{R}$;
target functional $\theta(\cdot)$ (one of $\tauf$, $\cf$,
$\Delta_{\RMST}(t^{*}, \cdot)$, etc.);
basis $\phi : \mathcal{X} \to \mathbb{R}^{k}$;
evaluation points $\{x_{i}\}_{i=1}^{n_{\mathrm{ev}}}$;
flag \texttt{use\_bb} for Bayesian bootstrap.
\Statex
\State Assemble the design matrix $\Phi \in \mathbb{R}^{n_{\mathrm{ev}} \times k}$
with rows $\phi(x_{i})^{\top}$; centre and (optionally) scale columns of
$\Phi$.
\State Precompute the QR decomposition of $\Phi$ if \texttt{use\_bb} is
false; otherwise the linear system is re-solved at each iteration.
\For{$r = 1, \ldots, R$}
  \State Evaluate the target at all evaluation points:
         $\theta^{(r)}_{i} \gets \theta(x_{i}; \theta_{\text{model}}^{(r)})$
         for $i = 1, \ldots, n_{\mathrm{ev}}$.
  \If{\texttt{use\_bb}}
    \State Draw $w^{(r)} \sim \mathrm{Dirichlet}(1, \ldots, 1)$.
    \State $\gamma^{(r)} \gets (\Phi^{\top} W^{(r)} \Phi)^{-1}
                                \Phi^{\top} W^{(r)} \theta^{(r)}$.
  \Else
    \State $\gamma^{(r)} \gets \Phi^{+} \theta^{(r)}$ via the
           precomputed QR factorisation, where $\Phi^{+}$ denotes the
           Moore--Penrose pseudoinverse.
  \EndIf
  \State Store $\gamma^{(r)}$.
\EndFor
\State \Return $\{\gamma^{(r)}\}_{r=1}^{R}$.
\end{algorithmic}
\end{algorithm}

Posterior summaries are then computed in the obvious way:
\[
  \bar{\gamma}_{k} = \frac{1}{R}\sum_{r=1}^{R}\gamma_{k}^{(r)},
  \qquad
  \widehat{\mathrm{CrI}}_{\alpha}(\gamma_{k}) =
  \bigl[Q_{\alpha/2}(\gamma_{k}^{(r)}),\ Q_{1-\alpha/2}(\gamma_{k}^{(r)})\bigr],
  \qquad
  \widehat{\Pr}(\gamma_{k} > 0) = \frac{1}{R}\sum_{r=1}^{R}
  \mathbf{1}\{\gamma_{k}^{(r)} > 0\},
\]
with joint summaries obtained from the empirical joint distribution of
$\{\gamma^{(r)}\}$ (e.g.\ posterior probability that $\gamma_{j}$ and
$\gamma_{k}$ have the same sign, or that a chosen contrast exceeds a
threshold).

\subsection{Implementation notes}

\paragraph{Basis selection.} The basis $\phi$ should be small and fixed:
five to ten terms is typically enough for an interpretable summary, and a
fixed basis is essential for validity of the projection-posterior
construction. If basis selection is desired, do it once on a point summary
of the posterior (e.g.\ run a lasso of the posterior mean of $\theta(\cdot)$
on a candidate dictionary), and then fix the selected basis for the
projection step. Repeating selection at every MCMC iteration would
introduce non-deterministic dependence on the draws and break the clean
``pushforward'' interpretation of \citet{Woody2020}; this is feasible but
adds complexity that we do not pursue.

\paragraph{Centring and scaling.} Centre and scale columns of $\Phi$ before
the projection so that coefficient magnitudes are comparable and the
intercept absorbs the mean of $\theta^{(r)}$. Standard regression hygiene;
matters more once relative magnitudes of $\gamma_{k}$ are interpreted
visually.

\paragraph{Evaluation set.} The choice of evaluation set determines the
implicit target population of the summary. Pooled RCT$+$OS rows describe
the combined sample; OS rows alone describe the target observational
population (the natural choice under (G1)); RCT rows describe the trial
population (typically the natural choice for HTA-style extrapolation under
(G2)). Always state explicitly which set is used.

\paragraph{Tree projection (CART summaries).} The procedure generalises
to projection onto a regression tree by replacing the OLS step with a
tree fit. The complication is that the tree \emph{structure} can change
across MCMC draws, which makes summarisation awkward. Two clean options:
(i) fix the tree structure (e.g.\ the tree fit to the posterior mean of
$\theta$) and refit only the leaf values at each iteration --- this stays
within the Woody--Carvalho--Murray framework; (ii) report subgroup
membership probabilities, summarising the variability of the partition
itself. Option (i) is simpler and is recommended unless the partition
itself is the object of interest.

\paragraph{Choice of $t^{*}$ for the RMST projection.} If $t^{*} \leq
t_{\RCT}$ the projection summary reflects RCT-anchored information; if
$t^{*} > t_{\RCT}$ the summary is computed from the extrapolated posterior
and depends on the HDP and the OS. Reporting projections at both an
in-window and an out-of-window $t^{*}$ traces how the heterogeneity story
evolves into the extrapolation regime, complementing the
coverage-vs-horizon diagnostic from the simulation.

\paragraph{Caveat on interpretation.} The projection is a \emph{summary},
not the truth. If $\tauf(\X)$ is genuinely nonlinear in covariates, a
linear projection captures the best linear approximation in MSE; the
nonlinear structure is absorbed into residual variation and not reported.
This is the price of interpretability and should be acknowledged when the
summary is presented.

\paragraph{Computational cost.} Negligible. Each iteration costs one
$k$-dimensional linear solve on $\Phi^{\top}\Phi$ (or a single
back-substitution with a cached QR), plus the cost of evaluating
$\theta^{(r)}(x_{i})$ for each $i$. The latter dominates only when
$\theta = \Delta_{\RMST}$, which requires a numerical integration per
evaluation point; in that case the integration can be cached at fixed
quadrature nodes across iterations.

\section{Related work}

The reframing places the project at the intersection of three literatures
that have historically not communicated very much: causal data fusion,
Bayesian survival extrapolation for HTA, and flexible Bayesian survival
modelling. We sketch each below and indicate where the proposed method sits.

\subsection{Survival extrapolation for HTA}

The foundational reference is NICE DSU Technical Support Document~14
\citep{Latimer2013}, which codified the practice of fitting standard
parametric AFT/PH families (exponential, Weibull, log-logistic, lognormal,
Gompertz, generalised gamma) to trial data and selecting among them by
fit and clinical plausibility. The TSD~14 process is widely followed but
also widely criticised: fits in the observed window are nearly
indistinguishable, while extrapolations diverge dramatically
\citep{BellGorrod2019, Bagust2014}.

External data have been used in HTA extrapolation in several distinct ways:
\begin{itemize}[leftmargin=*]
  \item \textbf{General-population mortality} as a lower bound on the
        hazard via additive-hazards (relative survival) models;
  \item \textbf{Registry conditional survival} from cancer registry
        databases used as external survival information for cancer
        \citep{Guyot2017};
  \item \textbf{Expert elicitation} treated as prior data on long-term hazard
        \citep{Guyot2017};
  \item \textbf{Longer-follow-up surrogate trials} such as the
        cilta-cel/LEGEND-2 example \citep{Vyas2023}, in which a phase~I
        study with longer follow-up informs extrapolation of a pivotal
        phase~II trial through informative priors on shape parameters.
\end{itemize}

The most directly comparable methodological piece is \citet{Jackson2023}'s
\texttt{survextrap}, which fits an M-spline hazard with a hierarchical prior
on the coefficients, jointly to individual-level RCT data and external
(possibly aggregate) data using Bayesian multi-parameter evidence synthesis.
\texttt{survextrap} supports proportional and flexible non-proportional
hazards, waning treatment effects, cure models, and additive-hazards
formulations \citep{Jackson2023}. Empirical evaluations
\citep{Timmins2025} report that the model gives accurate long-term
extrapolations when long-term external data on the patient population of
interest are included, and quantifies structural uncertainty appropriately
when they are not.

\paragraph{Where the proposed method differs from \texttt{survextrap}.}
\texttt{survextrap} does not address unmeasured confounding in the external
data; the external sources it accommodates are general-population mortality,
registry conditional survival, and elicited beliefs --- not a confounded
observational study of the \emph{same} treatment. The proposed framework
adds an explicit confounding-function correction $\cf(\X)$, allowing the
external data to be a contemporaneous OS subject to non-randomised
treatment assignment, not merely a background hazard reference. The
modelling style also differs: M-spline hazards versus BART forests on
AFT components, and a flexible parametric error versus an HDP nonparametric
error.

\subsection{Causal data fusion under unmeasured confounding}

The methodological core --- a confounding function $\cf(\X)$ correcting OS
contrasts that do not hold unconfoundedness --- is established in
\citet{YangLiuZengWang2025}, who derive semiparametric efficiency bounds for
joint estimation of $\tauf$ and $\cf$. \citet{Kallus2018} propose the
``experimental grounding'' framework, in which an RCT supplies a bias
function correcting an observational HTE estimator. \citet{Yao2024} extend
data fusion to partial OS treatment information. \citet{vanderLaan2024}
develop adaptive TMLE for RCT--RWD fusion that adapts borrowing to detected
bias.

The closest match to the proposed work is \citet{Ye2025}, who develop an
integrative AFT framework for high-dimensional RCT and RWD with censoring
and hidden confounding, using Stute-weighted penalised regression with
separate regularisation for $\tauf$ and the confounding component. The
proposed method is essentially a Bayesian-nonparametric counterpart to
\citet{Ye2025}: BART ensembles in place of penalised parametric models,
nonparametric error in place of a Gaussian or assumed AFT family, and
commensurate-prior structure on the baseline split between sources.

A practical caveat from \citet{Lin2024}: bias-correction methods for
marginal ATE estimation in this literature have not delivered
RMSE improvements over RCT-only analysis in their simulations, and the
methods' theoretical advantage is for the \emph{heterogeneous} effect
function rather than its average. The proposed framing targets HTE
plus long-term extrapolation, both of which are tasks for which RCT-only
inference is genuinely deficient --- the latter, simply because it cannot
be done from the trial alone.

\subsection{Bayesian flexible survival modelling}

The BART-based modelling of $\muf$, $\gf$, $\tauf$, $\cf$ rests on the
Bayesian causal forest (BCF) line of work \citep{Hahn2020}, which introduced
the prognostic/treatment-effect decomposition with separate prior scales
that the four-forest version generalises. The HDP error builds on the
centred stick-breaking mixture of \citet{Yang2010} and the wider
Bayesian-nonparametric survival literature; the hierarchical extension
specifically lets the source label index the group hierarchy, with shared
atoms across sources.

\subsection{Positioning summary}

\begin{table}[h]
\centering\small
\begin{tabular}{lccc}
\toprule
& Confounding & Long-term      & Nonparametric \\
& correction  & extrapolation  & error         \\
\midrule
\citet{Latimer2013} (TSD 14 workflow)        & --- & yes & ---  \\
\citet{Guyot2017}                            & --- & yes & ---  \\
\citet{Jackson2023} (\texttt{survextrap})    & --- & yes & flexible \\
\citet{Ye2025}                               & yes & limited$^\dagger$ & --- \\
\citet{YangLiuZengWang2025}                  & yes & ---       & --- \\
\citet{Kallus2018}, \citet{Yao2024}          & yes & ---       & --- \\
\textbf{This work}                           & \textbf{yes} & \textbf{yes} & \textbf{yes (HDP)} \\
\bottomrule
\end{tabular}
\caption{Schematic positioning. $\dagger$: \citet{Ye2025} targets HTE
estimation in the joint AFT framework; long-term extrapolation is not its
stated goal but is in principle compatible with the AFT structure.}
\end{table}

The niche is clear: the HTA-pragmatic literature handles extrapolation
flexibly but does not correct for confounding in the external source; the
causal-fusion literature corrects for confounding but does not target
extrapolation; and few methods combine the two with a fully nonparametric
error distribution. The proposed method sits deliberately in that gap.

\begin{thebibliography}{99}

\bibitem[Antonelli et al.(2019)]{Antonelli2019}
Antonelli, J., Parmigiani, G., and Dominici, F. (2019).
\newblock High-dimensional confounding adjustment using continuous spike
and slab priors.
\newblock \emph{Bayesian Analysis}, 14(3):805--828.

\bibitem[Bagust and Beale(2014)]{Bagust2014}
Bagust, A.\ and Beale, S. (2014).
\newblock Survival analysis and extrapolation modeling of time-to-event
clinical trial data for economic evaluation: an alternative approach.
\newblock \emph{Medical Decision Making}, 34(3):343--351.

\bibitem[Bell Gorrod et al.(2019)]{BellGorrod2019}
Bell Gorrod, H., Kearns, B., Stevens, J., Thokala, P., Labeit, A.,
Latimer, N., Tyas, D., and Sowdani, A. (2019).
\newblock A review of survival analysis methods used in NICE technology
appraisals of cancer treatments: consistency, limitations, and areas for
improvement.
\newblock \emph{Medical Decision Making}, 39(8):899--909.

\bibitem[Guyot et al.(2012)]{Guyot2012}
Guyot, P., Ades, A.~E., Ouwens, M.~J.~N.~M., and Welton, N.~J. (2012).
\newblock Enhanced secondary analysis of survival data: reconstructing the
data from published Kaplan--Meier survival curves.
\newblock \emph{BMC Medical Research Methodology}, 12:9.

\bibitem[Guyot et al.(2017)]{Guyot2017}
Guyot, P., Ades, A.~E., Beasley, M., Lueza, B., Pignon, J.-P., and
Welton, N.~J. (2017).
\newblock Extrapolation of survival curves from cancer trials using external
information.
\newblock \emph{Medical Decision Making}, 37(4):353--366.

\bibitem[Hahn et al.(2020)]{Hahn2020}
Hahn, P.~R., Murray, J.~S., and Carvalho, C.~M. (2020).
\newblock Bayesian regression tree models for causal inference:
regularization, confounding, and heterogeneous effects.
\newblock \emph{Bayesian Analysis}, 15(3):965--1056.

\bibitem[Jackson(2023)]{Jackson2023}
Jackson, C.~H. (2023).
\newblock survextrap: a package for flexible and transparent survival
extrapolation.
\newblock \emph{BMC Medical Research Methodology}, 23:282.

\bibitem[Kallus et al.(2018)]{Kallus2018}
Kallus, N., Puli, A.~M., and Shalit, U. (2018).
\newblock Removing hidden confounding by experimental grounding.
\newblock In \emph{Advances in Neural Information Processing Systems}, 31.

\bibitem[Latimer(2013)]{Latimer2013}
Latimer, N.~R. (2013).
\newblock Survival analysis for economic evaluations alongside clinical
trials---extrapolation with patient-level data: inconsistencies,
limitations, and a practical guide.
\newblock \emph{Medical Decision Making}, 33(6):743--754.
\newblock (NICE DSU Technical Support Document 14.)

\bibitem[Lin et al.(2024)]{Lin2024}
Lin, X., Tarp, J.~M., and Evans, R.~J. (2024).
\newblock Data fusion for efficiency gain in ATE estimation: a practical
review with simulations.
\newblock \emph{arXiv preprint} arXiv:2407.01186.

\bibitem[Rubin(1981)]{Rubin1981}
Rubin, D.~B. (1981).
\newblock The Bayesian bootstrap.
\newblock \emph{The Annals of Statistics}, 9(1):130--134.

\bibitem[Timmins et al.(2025)]{Timmins2025}
Timmins, I.~R., Jackson, C.~H., and colleagues (2025).
\newblock A simulation and case study to evaluate the extrapolation
performance of flexible Bayesian survival models when incorporating
real-world data.
\newblock \emph{arXiv preprint} arXiv:2505.16835.

\bibitem[van der Laan et al.(2024)]{vanderLaan2024}
van der Laan, M.~J., Qiu, S., and van der Laan, L. (2024).
\newblock Adaptive-TMLE for the average treatment effect based on
randomized controlled trial augmented with real-world data.
\newblock \emph{arXiv preprint} arXiv:2405.07186.

\bibitem[Vyas et al.(2023)]{Vyas2023}
Vyas, M., Hettle, R., Rana, M., Stevenson, K., et al. (2023).
\newblock Extrapolation of survival data using a Bayesian approach: a case
study leveraging external data from cilta-cel therapy in multiple myeloma.
\newblock \emph{Advances in Therapy}, 40(11):4877--4893.

\bibitem[Woody et al.(2020)]{Woody2020}
Woody, S., Carvalho, C.~M., and Murray, J.~S. (2020).
\newblock Bayesian inference with posterior projection.
\newblock \emph{Journal of Computational and Graphical Statistics},
29(4):798--808.

\bibitem[Yang et al.(2010)]{Yang2010}
Yang, M., Dunson, D.~B., and Baird, D. (2010).
\newblock Semiparametric Bayes hierarchical models with mean and variance
constraints.
\newblock \emph{Computational Statistics \& Data Analysis},
54(9):2172--2186.

\bibitem[Yang et al.(2025)]{YangLiuZengWang2025}
Yang, S., Liu, S., Zeng, D., and Wang, X. (2025).
\newblock Data fusion methods for the heterogeneity of treatment effect and
confounding function.
\newblock \emph{Bernoulli}, 31(4):2987--3012.

\bibitem[Yao et al.(2024)]{Yao2024}
Yao, D., Tang, C., Cui, Q., and Li, L. (2024).
\newblock Combining incomplete observational and randomized data for
heterogeneous treatment effects.
\newblock In \emph{Proceedings of the 33rd ACM International Conference
on Information and Knowledge Management}.

\bibitem[Ye et al.(2025)]{Ye2025}
Ye, X., Yang, S., Wang, X., et al. (2025).
\newblock Integrative analysis of high-dimensional RCT and RWD subject to
censoring and hidden confounding.
\newblock \emph{Lifetime Data Analysis}, 31:473--497.

\end{thebibliography}

\end{document}